\section{Input Contracts and Data Quality}

\subsection{What one count vector means}

A count vector is useful only when training and inference give its entries the
same meaning. The current model therefore stores an input contract containing
the device identity, MAC address, ordered categories, trained duration, and
counting-policy identifier. The policy is
\path{device-mac-bidirectional-ssdp-exclusive-v1}: select either Ethernet
direction and matching ARP hardware addresses, count TCP and ARP directly,
classify UDP source/destination port 1900 as SSDP, and exclude it from ordinary
UDP. Unsupported protocols are ignored by the four-category representation.

MAC identity selects traffic from mixed-device captures; it is not inferred
from an IP address that might change. The device name is the human-readable
identity in events and artefacts. The input contract binds both identities and
the counting interpretation to the model fingerprint.

The edge producer supplies completed counts rather than raw packets. The
consumer can check declared identity, dimensions, duration, timestamps, and
values; it cannot independently prove that an external producer counted every
packet correctly. That producer must implement the documented contract.

\subsection{Values, silence, and missing observations}

\path{WindowRecord} rejects negative counts, booleans, numeric strings,
fractional values, non-finite values, mismatched category lengths, and invalid
time bounds. Each count and the total must not exceed $2^{63}-2$ so the
implementation can add one without overflowing its integer range.
This is an input-range check, not a promise of exact floating-point arithmetic
for arbitrarily large counts.

An observed window contains traffic; an observed-silent window contains zeros;
a missing window contains no count vector. Missing data are not scored or
learned from. A silent vector is a legitimate observation and receives the
model's probability for $N=0$. Its raw score is zero, so a negative lower-tail
threshold will not flag silence simply because no packets were seen.
Detecting an outage from silence would need an additional, explicitly defined
signal; it is not supplied by this conditional composition model.

Training, calibration, and shared scoring check the configured duration and
category order. Dataset metadata must match the selected capture and device.
Streaming additionally enforces the UTC grid and chronological non-overlap.
The batch metadata checks do not replace the separate temporal branch's
complete cross-capture timeline validation.

\subsection{Inventory, structural validation, and fingerprints}

These are distinct gates. Inventory establishes whether configured files exist,
their size and checksum, and whether identical source bytes require an explicit
duplicate explanation. Structural validation reads classic PCAP to physical
EOF and distinguishes recoverable warnings from fatal corruption.
Fingerprinting identifies source bytes together with the policy that shapes
them into data. Passing one gate does not imply the others have been run.

The classic-PCAP validator checks global and packet headers, byte order,
microsecond/nanosecond resolution, version, snapshot length, packet boundaries,
and captured/original length relationships. Out-of-order timestamps,
snapshot-truncated packets, and zero-length records are counted as warnings,
with at most ten examples per warning code. Incomplete records, impossible
lengths, missing sources, or unsupported formats fail after a diagnostic report
has been written. PCAPNG is not supported by this structural validator, even
though the packet-reading path can encounter Scapy-supported PCAPNG sources.

The retained D-Link validation report records 12 configured captures, 721,934
packets, no fatal error, and 26 warnings in one capture. All twelve were parsed
to EOF. These are stored structural results, not a fresh re-read of every raw
capture for this report. UNSW preparation must not be described as having the
same complete classic-PCAP validation evidence when its source format differs.

A dataset hash covers ordered source hashes and sizes, category order, timestamp
policy, half-open windowing, duration, and partitions. Absolute clone paths and
generation timestamps are excluded. A model fingerprint instead covers the
parameters and inference semantics. A snapshot identifier names one accepted
version. These identifiers solve different reproducibility problems.

\subsection{Format changes and historical records}

Models and detector snapshots now use version 2 with the embedded input
contract. Older models without it fail clearly and must be regenerated with
matching thresholds; changing a configuration cannot prove how an old model was
trained. The bundled model pairs were regenerated without changing their
numeric alpha vectors or thresholds. Historical experiment records are retained
as records of their original revisions. A newer fingerprint need not have the
same bytes as an older fingerprint even when the numerical decisions agree.
